\chapter{Contexto Institucional e Alcance do Manual}
\label{chap:contexto}

A Universidade do Distrito Federal Professor Jorge Amaury Maia Nunes -- UnDF é uma fundação pública integrante da administração indireta do Distrito Federal, com autonomia acadêmica, administrativa e financeira nos termos de sua legislação. A Universidade oferece formação superior pública e gratuita e articula ensino, pesquisa, extensão, inovação e relação com a comunidade em diferentes áreas do conhecimento.

A relação atualizada de cursos, unidades acadêmicas, modalidades e regulamentos específicos deve ser consultada no portal e nos atos oficiais da UnDF. O manual não fixa número de cursos nem substitui informações acadêmicas sujeitas a atualização.

\section{Base institucional}

Este manual deve ser interpretado em conjunto com a Lei Complementar distrital nº 987/2021, o Decreto distrital nº 42.333/2021, o Estatuto, o Regimento Geral, o Regimento da Graduação aprovado pela Instrução nº 24/2025, a Resolução nº 12/2023 sobre elaboração e reestruturação de PPCs e os atos posteriores aplicáveis.

\section{Regras gerais e regras de curso}

Este documento apresenta orientações gerais. Cada PPC ou regulamento pode definir modalidades admitidas, carga horária, pré-requisitos, condições de orientação, composição da banca, critérios de avaliação, forma de defesa, prazos, entrega e depósito. Em caso de conflito, deve-se consultar a coordenação e a instância institucional competente.

\regracurso

\section{Independência da ferramenta de edição}

O presente documento explica o percurso do TCC, a estrutura do trabalho, a escrita, a metodologia e a normalização sem exigir conhecimento de programas de edição.

As orientações acadêmicas permanecem válidas independentemente da ferramenta autorizada para redigir e diagramar o trabalho. Recursos automáticos podem auxiliar a aplicação do padrão, mas não alteram os requisitos acadêmicos, a responsabilidade autoral nem as regras do curso.

\fluxopendente

Os materiais editáveis de apoio devem ser publicados em canal institucional controlado, enquanto este manual permanece centrado nas orientações acadêmicas. A localização, a unidade responsável e o fluxo de manutenção ainda devem ser confirmados.

\section{Como reconhecer o tipo de orientação}

Ao longo do manual, as orientações são identificadas conforme sua origem: requisito normativo; recomendação normativa; decisão institucional validada no âmbito do projeto; regra específica de curso ou PPC; boa prática acadêmica ou editorial; e fluxo institucional ainda pendente de definição. Essa distinção evita atribuir à ABNT escolhas da UnDF ou de cada curso e impede que decisões já consolidadas sejam tratadas novamente como propostas.
